\documentclass[11pt,a4paper]{article}
\usepackage[utf8]{inputenc}
\usepackage[T1]{fontenc}
\usepackage{amsmath,amssymb,amsthm}
\usepackage{listings}
\usepackage{geometry}
\usepackage{hyperref}
\geometry{margin=2.5cm}

\newtheorem{theorem}{Theorem}
\newtheorem{lemma}[theorem]{Lemma}

\lstset{
  language=C++,
  basicstyle=\ttfamily\small,
  keywordstyle=\bfseries,
  breaklines=true,
  numbers=left,
  numberstyle=\tiny,
  frame=single,
  tabsize=2
}

\title{IOI 2022 --- Radio Towers}
\author{}
\date{}

\begin{document}
\maketitle

\section{Problem Statement}

There are $N$ towers at positions $0, 1, \ldots, N{-}1$ with distinct
heights $H[0], \ldots, H[N{-}1]$.  Towers $i < j$ can
\emph{communicate} with parameter $D$ if there exists $k$ with
$i < k < j$ and $H[k] \ge \max(H[i], H[j]) + D$.

Given queries $(L, R, D)$, find the maximum number of towers in $[L, R]$
such that every pair can communicate.

\section{Solution}

\subsection{Characterization}

\begin{lemma}\label{lem:consec}
A set $S = \{s_1 < s_2 < \cdots < s_m\} \subseteq [L,R]$ is mutually
communicating with parameter $D$ if and only if every \emph{consecutive}
pair $(s_i, s_{i+1})$ can communicate.
\end{lemma}

\begin{proof}
The ``only if'' direction is trivial.  For ``if'': take any $s_a < s_b$
in $S$.  For each consecutive pair $(s_i, s_{i+1})$ there is a peak
$p_i$ with $H[p_i] \ge \max(H[s_i], H[s_{i+1}]) + D$.  Let
$p^* = \arg\max_i H[p_i]$ among the peaks for pairs between $s_a$ and
$s_b$.  Then $H[p^*] \ge \max(H[s_a], H[s_b]) + D$ because the maximum
among intermediate peaks dominates both endpoints.
\end{proof}

By Lemma~\ref{lem:consec}, the problem reduces to: find the longest
subsequence of $[L,R]$ such that between every two consecutive selected
towers there is a peak exceeding both by at least $D$.

\subsection{Greedy algorithm}

Scan left to right.  Greedily select a tower if:
\begin{enumerate}
  \item There exists a peak of sufficient height between it and the
        previously selected tower, \textbf{and}
  \item its height is low enough to be ``reachable'' from the next peak.
\end{enumerate}

More precisely: maintain the last selected tower.  For a candidate tower~$i$,
check (using a sparse table for range-max queries) whether
$\max_{k \in (\mathrm{last}, i)} H[k] \ge \max(H[\mathrm{last}], H[i]) + D$.
If so, select~$i$.  If not, and $H[i] < H[\mathrm{last}]$, replace the
last selection with~$i$ (a lower tower is strictly more flexible as a
chain endpoint).

\subsection{Correctness}

\begin{theorem}
The greedy algorithm produces an optimal set.
\end{theorem}
\begin{proof}[Proof sketch]
Suppose the greedy selects towers $g_1, g_2, \ldots, g_m$ and an optimal
solution selects $o_1, o_2, \ldots, o_k$ with $k > m$.  By an exchange
argument, we can show $H[g_i] \le H[o_i]$ for all $i$: if the greedy
tower is no taller, it is at least as good a ``continuation point.''
The greedy never misses a valid extension since it always holds the
lowest feasible endpoint, so $k \le m$, a contradiction.
\end{proof}

\section{Implementation}

\begin{lstlisting}
#include <bits/stdc++.h>
using namespace std;

int N;
vector<int> H;
vector<vector<int>> sparse;
int LOG;

void init(int _N, vector<int> _H) {
    N = _N;
    H = _H;
    LOG = 1;
    while ((1 << LOG) <= N) LOG++;
    sparse.assign(LOG, vector<int>(N));
    for (int i = 0; i < N; i++) sparse[0][i] = i;
    for (int k = 1; k < LOG; k++)
        for (int i = 0; i + (1 << k) <= N; i++) {
            int a = sparse[k-1][i];
            int b = sparse[k-1][i + (1 << (k-1))];
            sparse[k][i] = (H[a] >= H[b]) ? a : b;
        }
}

int rmq(int l, int r) { // index of max in [l, r]
    int k = __lg(r - l + 1);
    int a = sparse[k][l], b = sparse[k][r - (1 << k) + 1];
    return (H[a] >= H[b]) ? a : b;
}

int max_towers(int L, int R, int D) {
    if (L == R) return 1;

    int count = 0;
    int last = -1; // index of last selected tower

    for (int i = L; i <= R; i++) {
        if (last == -1) {
            // First tower: pick the first feasible starting point
            last = i;
            count = 1;
        } else {
            // Check if we can extend by selecting tower i
            if (last + 1 <= i - 1) {
                int mx = rmq(last + 1, i - 1);
                if (H[mx] >= H[last] + D && H[mx] >= H[i] + D) {
                    last = i;
                    count++;
                    continue;
                }
            }
            // If cannot extend, try replacing last with a lower tower
            if (H[i] < H[last]) {
                if (count == 1) {
                    last = i; // replace the only selected tower
                }
            }
        }
    }
    return max(count, 1);
}
\end{lstlisting}

\section{Complexity}

\begin{itemize}
  \item \textbf{Preprocessing:} $O(N \log N)$ for the sparse table.
  \item \textbf{Per query:} $O(R - L)$ --- each tower is visited once;
        range-max queries take $O(1)$.
  \item \textbf{Space:} $O(N \log N)$.
\end{itemize}

For full marks, a more sophisticated approach using the Cartesian tree
structure and offline/online segment-tree techniques can achieve
$O(\log N)$ or $O(\sqrt{N} \log N)$ per query.

\end{document}
